\documentclass[10pt,a4paper,landscape]{article}

\usepackage[margin=10mm]{geometry}
\usepackage{fontspec}
\usepackage{xcolor}
\usepackage{array}
\usepackage{booktabs}
\usepackage{tabularx}
\usepackage{longtable}
\usepackage{enumitem}
\usepackage{multicol}
\usepackage[most]{tcolorbox}
\usepackage{titlesec}
\usepackage{fancyhdr}
\usepackage{hyperref}
\usepackage{textcomp}
\usepackage{ragged2e}

\IfFontExistsTF{Arial}
  {\setmainfont{Arial}\setsansfont{Arial}}
  {\IfFontExistsTF{Calibri}
    {\setmainfont{Calibri}\setsansfont{Calibri}}
    {\IfFontExistsTF{Aptos}
      {\setmainfont{Aptos}\setsansfont{Aptos}}
      {\setmainfont{Latin Modern Sans}\setsansfont{Latin Modern Sans}}}}
\renewcommand{\familydefault}{\sfdefault}

\definecolor{Ink}{HTML}{1F2933}
\definecolor{Muted}{HTML}{586575}
\definecolor{Line}{HTML}{D6DEE6}
\definecolor{Panel}{HTML}{F5F7FA}
\definecolor{Blue}{HTML}{1E5AA8}
\definecolor{Green}{HTML}{237A57}
\definecolor{Amber}{HTML}{9A6A00}
\definecolor{Red}{HTML}{B42318}
\definecolor{Dark}{HTML}{111827}

\hypersetup{
  colorlinks=true,
  linkcolor=Blue,
  urlcolor=Blue
}

\setlength{\parindent}{0pt}
\setlength{\parskip}{4pt}
\setlist[itemize]{leftmargin=*, topsep=2pt, itemsep=2pt}
\setlist[enumerate]{leftmargin=*, topsep=2pt, itemsep=2pt}
\renewcommand{\arraystretch}{1.25}

\pagestyle{fancy}
\fancyhf{}
\lhead{\small TinyML Smart Plug User Manual}
\rhead{\small A4 Landscape Booklet Draft}
\cfoot{\small \thepage}
\renewcommand{\headrulewidth}{0.3pt}
\renewcommand{\footrulewidth}{0pt}

\titleformat{\section}
  {\Large\bfseries\color{Dark}}
  {}
  {0pt}
  {}
\titleformat{\subsection}
  {\bfseries\color{Dark}}
  {}
  {0pt}
  {}

\newcommand{\code}[1]{\texttt{\detokenize{#1}}}
\newcommand{\degreeC}{\textdegree C}
\newcommand{\Spread}[2]{%
  \clearpage
  \phantomsection
  \addcontentsline{toc}{section}{#1 #2}
  {\large\color{Muted}#1}\par
  {\Huge\bfseries\color{Dark}#2}\par
  \vspace{3mm}
}

\tcbset{
  sharp corners,
  boxrule=0.5pt,
  left=2mm,
  right=2mm,
  top=1.5mm,
  bottom=1.5mm,
  before skip=2mm,
  after skip=2mm
}

\newtcolorbox{warningbox}[1][]{
  colback=Red!5,
  colframe=Red!65!black,
  title=\textbf{Warning},
  #1
}

\newtcolorbox{notebox}[1][]{
  colback=Blue!5,
  colframe=Blue!65!black,
  title=\textbf{Note},
  #1
}

\newtcolorbox{layoutbox}[1][]{
  colback=Panel,
  colframe=Line,
  title=\textbf{Suggested Canva Layout},
  #1
}

\newtcolorbox{calloutbox}[1][]{
  colback=Green!5,
  colframe=Green!65!black,
  title=\textbf{Callout},
  #1
}

\begin{document}

\begin{titlepage}
\thispagestyle{empty}
\vspace*{12mm}
{\fontsize{34}{38}\selectfont\bfseries TinyML Smart Plug\par}
\vspace{5mm}
{\fontsize{18}{22}\selectfont Local socket-level protection for overload, socket heating, and series arc-fault events\par}
\vspace{16mm}
\begin{tcolorbox}[colback=Panel,colframe=Line,width=0.58\textwidth]
\textbf{User Manual Booklet}\\
Prototype operating guide for setup, monitoring, data logging, CSV review, OTA handling, and TinyML training workflows.
\end{tcolorbox}
\vfill
{\large Thesis Project Prototype\par}
{\large Firmware reference: \code{v7.8.2-p-gen0-final}\par}
\vspace{8mm}
\begin{layoutbox}
Use a full-page product photo or clean device render. Keep the title large and plain. Add only one short purpose line under the project name.
\end{layoutbox}
\end{titlepage}

\Spread{Spread 01}{Table of Contents}
\begin{multicols}{2}
\tableofcontents
\end{multicols}

\begin{notebox}
This booklet is written for a general user first. Admin, OTA, Logger, CSV Editor, and TinyML Trainer pages are for the project developer or trained operator.
\end{notebox}

\begin{layoutbox}
Two-column table of contents. Add small icons for safety, Wi-Fi, relay, data, firmware, and training if desired.
\end{layoutbox}

\Spread{Spread 02}{About the Device}
\begin{multicols}{2}
The TinyML Smart Plug is a smart plug prototype designed to add local protection at the socket level. It monitors voltage, current, estimated socket temperature, and waveform behavior to respond to overload, socket heating, and series arc-fault events.

The device can work with the PWA dashboard for monitoring, relay control, fault clearing, event history, data collection, CSV review, OTA controls, and TinyML development workflows.

Protection decisions are handled locally by the device firmware. The website improves visibility and control, but it is not required for the main protection loop.

\columnbreak

\textbf{What it helps detect}
\begin{itemize}
  \item Overload warning and sustained overload trip.
  \item Estimated socket heating warning and trip.
  \item Series arc-fault events using waveform features and TinyML inference.
  \item Abnormal voltage conditions, including under-voltage and over-voltage protection.
\end{itemize}

\textbf{What it is not}
\begin{itemize}
  \item Not a certified replacement for circuit breakers, AFCIs, GFCIs, proper wiring, or electrical inspection.
  \item Not a waterproof or outdoor-rated product.
  \item Not a commercial production unit.
\end{itemize}
\end{multicols}

\begin{warningbox}
Use this prototype only in safe, supervised test conditions. Do not use it with damaged outlets, burned plugs, unstable wiring, wet locations, or loads beyond the intended test range.
\end{warningbox}

\begin{layoutbox}
Left side: short description. Right side: simple block diagram from wall outlet to smart plug to protected appliance. Icons: plug, sensor, relay, display, Wi-Fi, cloud.
\end{layoutbox}

\Spread{Spread 03}{Key Features}
\begin{multicols}{2}
\textbf{Protection}
\begin{itemize}
  \item Local relay shutdown during critical fault conditions.
  \item Overload warning above 10.0 A.
  \item Sustained overload trip using averaged current behavior.
  \item Estimated socket heating warning at about 39.0\degreeC{} and trip at about 40.0\degreeC{}.
  \item Series arc-fault detection using waveform features and TinyML model output.
\end{itemize}

\textbf{User operation}
\begin{itemize}
  \item PWA dashboard at \code{tinyml-smart-plug.web.app}.
  \item Relay ON/OFF control from the dashboard.
  \item Clear Fault command from the dashboard.
  \item Manual relay ON button for limited offline use.
  \item OLED display for local status and basic readings.
\end{itemize}

\columnbreak

\textbf{Developer and admin tools}
\begin{itemize}
  \item Firebase live data and event logging when Wi-Fi is available.
  \item Logger for controlled data collection sessions.
  \item CSV viewer, plotter, editor, arc marker tool, and segment labeling.
  \item OTA firmware controls and manual firmware download workflow.
  \item TinyML data preparation, model training, benchmarking, and exported firmware headers.
\end{itemize}
\end{multicols}

\begin{notebox}
Voltage protection is included because the firmware and manuscript define under-voltage and over-voltage states that can lock relay control during unsafe mains conditions.
\end{notebox}

\begin{layoutbox}
Use three clean feature groups: Protection, User Operation, Admin Tools. Use check icons or simple line icons. Avoid long paragraphs.
\end{layoutbox}

\Spread{Spread 04}{Important Safety Warnings}
\begin{multicols}{2}
\begin{warningbox}
Always inspect the outlet, plug, and appliance before use. Do not use the device on a loose, damaged, burned, wet, or unstable outlet.
\end{warningbox}

\begin{warningbox}
If the device trips, do not immediately turn the relay back ON. Inspect the connected appliance, socket, plug, wiring, and load condition first.
\end{warningbox}

\begin{warningbox}
The prototype weighs around 750 g. A weak wall socket may not support it securely. Prevent sagging and falling before energizing the load.
\end{warningbox}

\columnbreak

\textbf{Do not use with}
\begin{itemize}
  \item Loose wall sockets or protruding outlet adapters.
  \item Burned, cracked, wet, corroded, or unstable outlets.
  \item Damaged extension cords or overloaded power strips.
  \item Appliances with visible plug, cord, or casing damage.
  \item Loads that exceed the intended prototype test conditions.
\end{itemize}

\textbf{Stop use immediately if}
\begin{itemize}
  \item There is smoke, burning smell, sparking, or unusual heat.
  \item The enclosure, plug, socket, or cable becomes damaged.
  \item The device repeatedly trips after clearing a fault.
  \item The plug no longer fits tightly in the outlet.
\end{itemize}
\end{multicols}

\begin{layoutbox}
Use a red safety band with four warning cards. Add icons for heat, unstable socket, electric shock, and inspect before reset.
\end{layoutbox}

\Spread{Spread 05}{Physical Handling and Mounting}
\begin{multicols}{2}
The TinyML Smart Plug is heavier than a normal plug-in device because of its relay, heatsink, enclosure, and protection circuitry. The current prototype weighs around 750 g.

Because of this weight, the device may sag if it is plugged into a loose, protruding, or weak wall socket. The plug must remain tightly connected during operation.

\textbf{Recommended mounting}
\begin{enumerate}
  \item Use a firm wall outlet that holds the plug tightly.
  \item Position the device so it sits flat or nearly flush against the wall.
  \item Support the device if needed using a safe stopper, backing, or spacer.
  \item Keep the device stable before connecting the appliance load.
\end{enumerate}

\columnbreak

\textbf{Using an extension wire or power strip}
\begin{itemize}
  \item Lay the smart plug flat on the floor or on a stable surface.
  \item Do not let the device hang from the outlet, cord, or power strip.
  \item Keep the power strip away from water, dust buildup, and foot traffic.
\end{itemize}

\begin{warningbox}
Do not use loose, damaged, burned, cracked, or unstable outlets. A loose connection can create heat and may make the socket heating condition worse.
\end{warningbox}
\end{multicols}

\begin{layoutbox}
Show a correct/incorrect mounting comparison: flush supported wall mount, flat power strip setup, sagging unsupported outlet. Add callouts for plug fit, support surface, and protected socket.
\end{layoutbox}

\Spread{Spread 06}{System Components}
\begin{tabularx}{\textwidth}{>{\bfseries}p{0.22\textwidth}X}
\toprule
Component & Purpose \\
\midrule
Red main rocker switch & Powers the TinyML Smart Plug itself. This is different from the relay control for the appliance socket. \\
Protected outlet/socket & Output socket where the appliance or test load is connected. It only receives power when the relay output is ON. \\
Relay / solid-state relay path & Connects or disconnects power to the protected outlet during user control or protection trips. \\
ESP32-S3 controller & Runs Wi-Fi, local protection logic, sensor processing, TinyML inference, Firebase communication, and control handling. \\
Voltage sensor & Measures mains voltage for monitoring and abnormal voltage protection. \\
Current sensor and ADC & Measures load current and waveform behavior used by protection and TinyML features. \\
Thermistor & Measures internal/sensor temperature used to estimate socket heating behavior. \\
OLED display & Shows local status and limited live values such as voltage, current, estimated socket temperature, and apparent power. \\
Buzzer/chime & Gives audible feedback for setup, collection, faults, and command acknowledgment. \\
Manual relay ON button & Red circular non-tactile button beside the main switch. Used only to turn the relay ON during limited offline operation. \\
\bottomrule
\end{tabularx}

\begin{notebox}
This manual uses the user-facing term \textbf{relay output}. The technical prototype uses a relay/SSR switching path to control the protected socket.
\end{notebox}

\begin{layoutbox}
Use one device photo with numbered callouts. Suggested labels: Main Switch, Manual Relay ON Button, OLED, Protected Outlet, Plug Input, Enclosure, Vent/Heat Area if visible.
\end{layoutbox}

\Spread{Spread 07}{Quick Start}
\begin{multicols}{2}
\textbf{Basic setup}
\begin{enumerate}
  \item Plug the TinyML Smart Plug into a stable outlet.
  \item Turn the red main switch ON.
  \item Connect the device to a 2.4 GHz Wi-Fi network.
  \item Open \code{tinyml-smart-plug.web.app} in Google Chrome.
  \item Plug the appliance or test load into the protected outlet.
  \item Turn the relay ON from the PWA rocker control.
\end{enumerate}

\textbf{Important behavior}
\begin{itemize}
  \item The relay is OFF by default after startup.
  \item Relay ON gives power to the appliance socket.
  \item Relay OFF removes power from the appliance socket.
  \item The red main switch powers the smart plug itself.
\end{itemize}

\columnbreak

\textbf{If Wi-Fi is unavailable}
\begin{enumerate}
  \item Keep the device plugged in and powered.
  \item Connect the load first.
  \item Hold the manual relay ON button for about 5 seconds.
  \item Wait for the chime.
  \item Release the button.
\end{enumerate}

\begin{warningbox}
Use the manual relay ON button only when needed. It cannot turn the relay OFF. Use the PWA relay OFF command or remove power using the main switch if the load must be stopped.
\end{warningbox}
\end{multicols}

\begin{layoutbox}
Make this a visual flow: Power, Wi-Fi, PWA, Plug Load, Relay ON. Add a small offline branch for the manual relay ON button.
\end{layoutbox}

\Spread{Spread 08}{Wi-Fi Setup}
\begin{multicols}{2}
\textbf{First-time setup}
\begin{enumerate}
  \item Turn the red main switch ON.
  \item Open Wi-Fi settings on a phone or laptop.
  \item Connect to the \code{TinyML-Smart-Plug} access point.
  \item Wait for the captive portal to appear.
  \item Select \textbf{Configure Wi-Fi}.
  \item Choose a 2.4 GHz Wi-Fi network.
  \item Enter the Wi-Fi password.
  \item Wait for the device to chime and show Wi-Fi connected status on the OLED.
\end{enumerate}

\textbf{After setup}
\begin{itemize}
  \item Saved credentials are kept after reset.
  \item On the next startup, the device tries to reconnect automatically.
  \item The PWA can monitor and control the device when Wi-Fi and Firebase are reachable.
\end{itemize}

\columnbreak

\begin{warningbox}
The device accepts 2.4 GHz Wi-Fi only. A 5 GHz-only network will not appear or will not connect.
\end{warningbox}

\begin{notebox}
Firmware source sets the access point name to \code{TinyML-Smart-Plug}. If no saved Wi-Fi credentials exist, the boot AP waits about 15 seconds for a client. After a client connects and the portal opens, normal setup has about 45 seconds to complete.
\end{notebox}
\end{multicols}

\begin{layoutbox}
Use four screenshots: phone Wi-Fi list, captive portal, Configure Wi-Fi page, OLED Wi-Fi connected message. Add a 2.4 GHz badge.
\end{layoutbox}

\Spread{Spread 09}{Resetting the Device}
\begin{multicols}{2}
Resetting restarts the smart plug and clears temporary runtime state. It does not erase saved Wi-Fi credentials.

\textbf{How to reset}
\begin{enumerate}
  \item Turn the red main switch OFF.
  \item Wait at least 10 seconds.
  \item Turn the red main switch ON again.
  \item Wait for startup and Wi-Fi reconnect.
\end{enumerate}

\textbf{Reset clears}
\begin{itemize}
  \item RAM state.
  \item Temporary buffers.
  \item Runtime timers.
  \item Active fault latch state after full restart.
\end{itemize}

\columnbreak

\textbf{Reset does not clear}
\begin{itemize}
  \item Saved Wi-Fi credentials.
  \item Firmware version.
  \item Cloud history already uploaded to Firebase.
  \item CSV files already downloaded or saved in the PWA workflow.
\end{itemize}

\begin{notebox}
The 10-second wait helps the device fully power down and gives capacitors time to discharge before restarting.
\end{notebox}
\end{multicols}

\begin{warningbox}
If the device tripped because of a fault, inspect the load, plug, socket, and wiring before resetting or turning the relay back ON.
\end{warningbox}

\begin{layoutbox}
Use a simple OFF, wait 10 seconds, ON diagram. Add a small note that Wi-Fi credentials remain saved.
\end{layoutbox}

\Spread{Spread 10}{Offline Protection}
\begin{multicols}{2}
The PWA is important for monitoring and control, but the protection loop is local. The device can still respond to overload, socket heating, abnormal voltage, and arc-fault conditions even when the website is not open.

\textbf{When Wi-Fi is unavailable}
\begin{itemize}
  \item Local protection remains active.
  \item The relay can still shut OFF during critical fault conditions.
  \item The OLED shows limited live information.
  \item The buzzer can still give local feedback.
\end{itemize}

\columnbreak

\textbf{Limitations without Wi-Fi}
\begin{itemize}
  \item PWA relay control is unavailable.
  \item Firebase live data and event uploads may be delayed or missing.
  \item Data collection sessions cannot be fully managed from the website.
  \item CSV viewing, editing, and downloading are not available until the PWA has data.
\end{itemize}

\begin{notebox}
The OLED can show limited live values such as voltage, current, estimated socket temperature, and apparent power.
\end{notebox}
\end{multicols}

\begin{layoutbox}
Split the page into Online and Offline columns. Show protection logic inside the device, not inside the website.
\end{layoutbox}

\Spread{Spread 11}{Using the PWA}
\begin{multicols}{2}
\textbf{Open the dashboard}
\begin{enumerate}
  \item Open Google Chrome.
  \item Go to \code{tinyml-smart-plug.web.app}.
  \item Allow the page to load the dashboard.
  \item Install the PWA if Chrome prompts installation.
\end{enumerate}

\textbf{For iPhone or iPad}
\begin{itemize}
  \item Open the site in Safari.
  \item Use \textbf{Share} then \textbf{Add to Home Screen}.
  \item Open the installed icon from the Home Screen.
\end{itemize}

\columnbreak

\textbf{Main dashboard areas confirmed in code}
\begin{itemize}
  \item System Overview.
  \item Relay rocker control.
  \item Refresh and Clear Fault controls.
  \item Status cards for load family, protection, fault timing, firmware, and last event.
  \item Main readings for voltage, current, sensor temperature, estimated socket temperature, expected normal temperature, and apparent power.
  \item Waveform disturbance and conduction intermittence feature cards.
  \item History section with date range, download, and clear actions.
  \item Admin-only Logger, Plot Viewer, CSV tools, OTA controls, and TinyML-related workflow support.
\end{itemize}
\end{multicols}

\begin{layoutbox}
Use a dashboard screenshot with numbered labels. Keep the PWA URL visible as a small footer callout.
\end{layoutbox}

\Spread{Spread 12}{Main PWA Interface}
\begin{tabularx}{\textwidth}{>{\bfseries}p{0.23\textwidth}X}
\toprule
Area & What the user sees or does \\
\midrule
System Overview & Shows device state, relay state, fault status, and quick controls. \\
Relay rocker & Turns the protected appliance socket ON or OFF when controls are available. \\
Clear Fault & Sends a fault clear command to the device. Use only after inspection. \\
Main Readings & Shows voltage, current, sensor temperature, estimated socket temperature, expected normal temperature, and apparent power. \\
Status cards & Show load family, protection state, fault timing, firmware information, and last event. \\
Feature cards & Show waveform and conduction features used for analysis and TinyML behavior. \\
History & Shows alerts/events with range filters such as Today, Yesterday, Last 7 days, and Last 30 days. \\
Menu & Contains notification mode, User/Admin mode switch, Change WiFi, OTA guide, and related admin controls. \\
\bottomrule
\end{tabularx}

\begin{notebox}
The PWA source uses a top-right menu button. This manual describes it as the three-line menu for clarity.
\end{notebox}

\begin{layoutbox}
Use one full-width PWA screenshot. Add callouts: Relay, Clear Fault, Readings, Status Cards, History, Menu.
\end{layoutbox}

\Spread{Spread 13}{Relay Control}
\begin{multicols}{2}
The relay rocker at the top of the PWA controls the relay output. This controls electricity to the protected appliance socket.

\textbf{Relay states}
\begin{itemize}
  \item \textbf{Relay ON:} the connected appliance receives power.
  \item \textbf{Relay OFF:} the connected appliance is turned off.
  \item \textbf{Red main switch ON:} the smart plug itself is powered.
  \item \textbf{Red main switch OFF:} the whole smart plug is powered down.
\end{itemize}

\columnbreak

\textbf{How to turn the appliance ON}
\begin{enumerate}
  \item Confirm the outlet and device are stable.
  \item Plug the appliance into the protected outlet.
  \item Open the PWA.
  \item Move the relay rocker to ON.
  \item Confirm that the load receives power.
\end{enumerate}

\begin{notebox}
The relay is OFF by default after startup. This prevents the connected appliance from energizing automatically.
\end{notebox}
\end{multicols}

\begin{warningbox}
If the PWA blocks Relay ON, do not force operation. The device may be offline, disconnected, in startup, in OTA/config mode, or in a fault condition.
\end{warningbox}

\begin{layoutbox}
Show a simple distinction: main switch powers device; relay rocker powers protected socket. Use two separate arrows.
\end{layoutbox}

\Spread{Spread 14}{Fault Warnings and Critical Trips}
\begin{multicols}{2}
Fault behavior may be a warning or a critical trip.

\textbf{Warning}
\begin{itemize}
  \item The device detects a condition that needs attention.
  \item The relay may remain ON while the condition is watched.
  \item The PWA and OLED may show warning status.
\end{itemize}

\textbf{Critical trip}
\begin{itemize}
  \item The relay turns OFF as a system response.
  \item The connected appliance loses power.
  \item The user cannot normally turn the relay back ON until the fault is cleared.
  \item The device may chime, update Firebase/PWA status, and log an event if Wi-Fi is available.
\end{itemize}

\columnbreak

\textbf{Fault types used by firmware}
\begin{itemize}
  \item Overload warning.
  \item Sustained overload trip.
  \item Heating warning and heating trip.
  \item Arc-fault trip.
  \item Under-voltage and over-voltage lockout.
  \item Unplugged, disconnected, startup, OTA, Wi-Fi, or safe-mode states may also limit controls.
\end{itemize}

\begin{notebox}
The firmware gives priority to critical states such as heating trip, arc, sustained overload, and voltage lockout when deciding relay behavior.
\end{notebox}
\end{multicols}

\begin{layoutbox}
Use a two-level status diagram: Warning means watch and inspect; Critical Trip means relay OFF and inspect before clearing.
\end{layoutbox}

\Spread{Spread 15}{Clearing Faults}
\begin{multicols}{2}
Faults can be cleared in two ways.

\textbf{Option 1: PWA Clear Fault}
\begin{enumerate}
  \item Open the PWA dashboard.
  \item Inspect the displayed fault condition.
  \item Inspect the appliance, plug, socket, wiring, and load condition.
  \item Click \textbf{Clear Fault}.
  \item Wait for the device chime.
  \item Turn the relay ON only if the setup is safe.
\end{enumerate}

\columnbreak

\textbf{Option 2: Device reset}
\begin{enumerate}
  \item Turn the red main switch OFF.
  \item Wait at least 10 seconds.
  \item Turn the red main switch ON.
  \item Wait for startup.
  \item Turn the relay ON only after inspection.
\end{enumerate}

\begin{notebox}
The firmware acknowledges the Clear Fault command with a local chime.
\end{notebox}
\end{multicols}

\begin{warningbox}
Clearing a fault means the user acknowledges that a problem occurred. Do not clear a fault just to continue operation. Inspect and address the connected appliance, socket, plug, wiring, or load condition first.
\end{warningbox}

\begin{layoutbox}
Use an inspection checklist graphic before the Clear Fault button. Add a red ``Inspect First'' label.
\end{layoutbox}

\Spread{Spread 16}{Manual Relay ON Button}
\begin{multicols}{2}
The red circular non-tactile button beside the main rocker switch is the manual relay ON button.

It is intended for limited use when Wi-Fi or the PWA is unavailable. It can turn the relay ON, but it cannot turn the relay OFF.

\textbf{How to use}
\begin{enumerate}
  \item Connect the load to the protected outlet first.
  \item Make sure the device is stable and powered.
  \item Hold the manual relay ON button for about 5 seconds.
  \item If a load is detected, the relay turns ON and the device chimes.
  \item Release the button after the chime.
\end{enumerate}

\columnbreak

\textbf{If no load is detected}
\begin{itemize}
  \item Holding the button for 5 seconds will not turn the relay ON.
  \item Confirm that a load is connected.
  \item Confirm that the appliance switch is in a state that draws detectable current.
\end{itemize}

\textbf{To turn the relay OFF}
\begin{itemize}
  \item Use the PWA relay OFF command when available.
  \item A critical fault can also force the relay OFF.
  \item Turn the red main switch OFF if the device must be fully powered down.
\end{itemize}
\end{multicols}

\begin{warningbox}
This button is directly connected to the relay driver path and can bypass software restrictions. It must not be stuck, blocked, accidentally held, taped down, or compromised by any object.
\end{warningbox}

\begin{notebox}
Firmware confirms a 5-second current-based physical relay ON confirmation. The manuscript and hardware notes describe this as the red manual relay ON button beside the main switch.
\end{notebox}

\begin{layoutbox}
Use a close-up photo of the red circular button with a warning callout. Add a small timing bar: hold 5 seconds, wait for chime, release.
\end{layoutbox}

\Spread{Spread 17}{Changing Wi-Fi}
\begin{multicols}{2}
Use Change Wi-Fi when the router changes, credentials are wrong, or the device must connect to a different 2.4 GHz network.

\textbf{From the PWA}
\begin{enumerate}
  \item Open the PWA dashboard.
  \item Click the three-line menu at the top-right.
  \item Click \textbf{Change WiFi}.
  \item Wait for the device access point to open.
  \item Open phone or laptop Wi-Fi settings.
  \item Connect to the \code{TinyML-Smart-Plug} access point.
  \item Follow the same Wi-Fi setup process.
\end{enumerate}

\columnbreak

\textbf{If the device is not connected}
\begin{itemize}
  \item Reset the device using the red main switch.
  \item Watch for the access point during startup.
  \item Connect and configure the new 2.4 GHz Wi-Fi network.
\end{itemize}

\begin{notebox}
Changing Wi-Fi does not change the firmware or delete historical cloud records.
\end{notebox}
\end{multicols}

\begin{notebox}
For Change Wi-Fi, firmware opens a manual AP wait window of about 45 seconds and allows a longer manual portal session of about 4 minutes once the portal is active.
\end{notebox}

\begin{layoutbox}
Use three screenshots: menu open, Change WiFi button, Wi-Fi captive portal. Add a small 2.4 GHz-only callout.
\end{layoutbox}

\Spread{Spread 18}{Admin Mode}
\begin{multicols}{2}
Admin mode opens development and maintenance tools. It is intended for the project developer or trained operator.

\textbf{How to open Admin mode}
\begin{enumerate}
  \item Open the PWA dashboard.
  \item Click the top-right menu.
  \item Find the \textbf{Mode} switch.
  \item Change from \textbf{User} to \textbf{Admin}.
\end{enumerate}

\textbf{Admin-only areas confirmed in code}
\begin{itemize}
  \item OTA controls.
  \item Logger controls.
  \item Original and Processed session tabs.
  \item CSV upload.
  \item Plot Viewer and CSV editor.
  \item Save and download processed CSV tools.
\end{itemize}

\columnbreak

\textbf{Admin controls in the menu}
\begin{itemize}
  \item Notification mode: Silent or Alert.
  \item Mode: User or Admin.
  \item Change WiFi.
  \item Revert Firmware.
  \item OTA Guide.
\end{itemize}

\begin{warningbox}
Admin tools can change data collection, firmware update, and processed dataset outputs. Use them only when the test setup is stable and the operator understands the workflow.
\end{warningbox}
\end{multicols}

\begin{layoutbox}
Use a menu screenshot with the User/Admin switch highlighted. Add a lock or wrench icon for developer tools.
\end{layoutbox}

\Spread{Spread 19}{OTA Firmware Update}
\begin{multicols}{2}
OTA controls are for the developer. Normal users should not update firmware unless instructed.

\textbf{Current status}
\begin{itemize}
  \item HTTPS OTA is present in the firmware and PWA.
  \item The user notes say the HTTPS OTA path is currently not working properly.
  \item Only the developer has access to firmware version releases.
\end{itemize}

\textbf{When to check updates}
\begin{itemize}
  \item Check when the firmware version shown on the website differs from the version shown in the menu or firmware information area.
  \item Current firmware reference in source: \code{v7.8.2-p-gen0-final}.
\end{itemize}

\columnbreak

\textbf{Manual update workflow}
\begin{enumerate}
  \item In Admin OTA, click the green download arrow to download the \code{.bin} firmware file.
  \item Open the device AP using \textbf{Change WiFi}, or reset the device if it is not connected.
  \item Connect to the TinyML Smart Plug AP.
  \item Turn off mobile data.
  \item Exit the captive portal window.
  \item Open a real browser window.
  \item Go to \code{192.168.4.1}.
  \item Click \textbf{Update Firmware}.
  \item Upload the downloaded \code{.bin} file.
  \item Wait for upload, flashing, and automatic restart.
\end{enumerate}
\end{multicols}

\begin{warningbox}
Do not remove power during firmware upload or flashing. Use a stable outlet and wait for the device to restart.
\end{warningbox}

\begin{notebox}
The main firmware and PWA support HTTPS OTA through the published version and firmware URL fields. The repository also contains a browser upload OTA utility route at \code{/update}. For the manual local upload workflow, leave the captive portal and use a real browser at \code{192.168.4.1}.
\end{notebox}

\begin{layoutbox}
Use a firmware update timeline: Download BIN, Open AP, Real Browser, Upload, Flash, Restart. Add a warning icon at the flashing step.
\end{layoutbox}

\Spread{Spread 20}{Data Collection / Logger}
\begin{multicols}{2}
The Logger records controlled sessions for analysis and TinyML dataset preparation. This section is for Admin users.

\textbf{Before collecting}
\begin{itemize}
  \item Use a stable outlet and safe mounting.
  \item Confirm the load type and load name.
  \item Decide the phase: Start, Steady, or Arc.
  \item Keep the setup supervised.
\end{itemize}

\textbf{Logger fields confirmed in PWA}
\begin{itemize}
  \item Duration in seconds.
  \item Load type or family.
  \item Device/load name.
  \item Trial number.
  \item Division tag: start, steady, or arc.
  \item Notes.
  \item Trusted normal switch.
\end{itemize}

\columnbreak

\textbf{How to collect}
\begin{enumerate}
  \item Open Admin mode.
  \item Go to the Logger section.
  \item Enter collection duration in seconds.
  \item Select load type/family.
  \item Enter load name and trial number.
  \item Select start, steady, or arc phase.
  \item Click or switch \textbf{Collect}.
  \item Wait for the device chime.
  \item Wait until the Collect rocker turns off automatically.
  \item Wait for buffered upload to finish.
\end{enumerate}
\end{multicols}

\begin{notebox}
The CSV time column may continue counting temporarily after collection ends because buffered samples are still uploading. When counting stops, the CSV is ready for viewing or downloading.
\end{notebox}

\begin{layoutbox}
Use a form screenshot with callouts for Duration, Load Type, Trial, Phase, Collect switch, and session tabs.
\end{layoutbox}

\Spread{Spread 21}{CSV Plotter and Editor}
\begin{multicols}{2}
The Plot Viewer opens logged or uploaded CSV sessions for review, graphing, labeling, and download.

\textbf{Opening CSV data}
\begin{itemize}
  \item Use Original sessions to view raw logger uploads.
  \item Use Processed sessions to view saved edited outputs.
  \item Use Upload CSV to open a local CSV file.
  \item If cloud saving fails after upload, the file can still open in the local viewer.
\end{itemize}

\textbf{Graphing and viewing}
\begin{itemize}
  \item The viewer plots time-based CSV columns.
  \item Series can be searched, shown, hidden, set to default, set to all, or set to none.
  \item Visible series can be assigned to Y1 or Y2 axis.
  \item Controls include normalize, smooth, curve style, follow, zoom, play/pause, speed, and scrubber.
\end{itemize}

\columnbreak

\textbf{Data columns and features}
\begin{itemize}
  \item Voltage, current, temperature, and apparent power can be viewed when present.
  \item Waveform and feature columns can be plotted when present.
  \item The PWA recognizes arc labels using \code{label_arc}.
  \item CSV metadata can include family, device/load name, trial, notes, trusted normal, and division tag.
\end{itemize}

\textbf{Saving}
\begin{itemize}
  \item Download CSV saves the current edited CSV.
  \item Save stores processed CSV output to the processed session area when Firebase saving is available.
  \item Download All can download sessions in order.
\end{itemize}
\end{multicols}

\begin{notebox}
The exact plotted series depends on the CSV headers available in the selected session. Missing columns do not appear as graphable series.
\end{notebox}

\begin{layoutbox}
Use a large plot screenshot. Add labels for Series panel, Zoom, Smooth, Normalize, Play/Pause, Save, and Download CSV.
\end{layoutbox}

\Spread{Spread 22}{Arc Markers and Section Labels}
\begin{multicols}{2}
The CSV editor supports arc marker editing and waveform section labeling for TinyML dataset preparation.

\textbf{Arc marker editing}
\begin{itemize}
  \item Enable Show arcs to display arc labels.
  \item Use the arc rows field to enter row numbers.
  \item Single rows and ranges are supported, such as \code{10-14,25}.
  \item Click Add to mark rows as arc.
  \item Click Remove to clear arc markers.
  \item Edited arc labels are stored in \code{label_arc}.
\end{itemize}

\columnbreak

\textbf{Section labels}
\begin{itemize}
  \item Use Label to select Start, Steady, or Arc.
  \item Use Start and End row fields to define a section.
  \item Set Start and Set End can use the current viewer position.
  \item Add creates a labeled segment.
  \item Remove deletes a selected segment.
  \item Split downloads can create separate processed CSV files from labeled sections.
\end{itemize}

\textbf{How phases are used}
\begin{itemize}
  \item Start rows are used for load context training.
  \item Steady and Arc rows are used for arc training data.
  \item Arc rows should represent the event period clearly.
\end{itemize}
\end{multicols}

\begin{warningbox}
Incorrect arc labels can reduce model quality. Review the waveform and event timing before saving processed CSV files.
\end{warningbox}

\begin{layoutbox}
Show the editor panel with the arc row input and segment list. Add callouts for Label, Start row, End row, Add, Remove, and Download Split.
\end{layoutbox}

\Spread{Spread 23}{TinyML Trainer Workflow}
\begin{multicols}{2}
The TinyML Trainer workflow is for the developer. It converts collected CSV logs into training datasets, model reports, and firmware header files.

\textbf{Folder workflow confirmed in code}
\begin{itemize}
  \item Raw CSV files go in \code{tinyml/data/{gen}/raw}.
  \item Generation folders use names such as \code{gen0}, \code{gen1}, and later generations.
  \item Prepared outputs are written to \code{tinyml/data/{gen}/arc_training.csv} and \code{tinyml/data/{gen}/load_context.csv}.
  \item Reports are written to \code{tinyml/benchmark/{gen}}.
  \item Active firmware headers are written to \code{tinyml/model}.
  \item Archive copies are kept in generation-specific archive folders.
\end{itemize}

\columnbreak

\textbf{Preparation step}
\begin{itemize}
  \item Run \code{prepare_data.py} with the target generation.
  \item The script reads raw CSV files from the generation raw folder.
  \item It cleans and normalizes CSV columns.
  \item Start rows are kept for context training.
  \item Steady and arc rows are kept for arc training.
  \item The output files are \code{arc_training.csv} and \code{load_context.csv}.
\end{itemize}

\textbf{Trainer GUI}
\begin{itemize}
  \item \code{tinyml/tinyml_trainer.py} provides a GUI workflow.
  \item The GUI shows the generation workflow and can help run preparation, training, and export tasks.
\end{itemize}
\end{multicols}

\begin{layoutbox}
Use a folder pipeline diagram: Raw CSVs, Prepare Data, Arc Dataset, Context Dataset, Train, Benchmark, Export Headers, Firmware.
\end{layoutbox}

\Spread{Spread 24}{TinyML Training and Export}
\begin{multicols}{2}
\textbf{Arc model training}
\begin{itemize}
  \item \code{train_rf.py} trains the Random Forest arc model.
  \item \code{train_et.py} trains the Extra Trees arc model.
  \item \code{train_duel.py} can compare RF and ET workflows.
  \item \code{train_feature_subsets.py} supports feature subset sweeps.
  \item Reports and feature summaries are saved in the benchmark folder.
\end{itemize}

\textbf{Context model training}
\begin{itemize}
  \item \code{train_context.py} trains the load family context model.
  \item The context model classifies six load families.
  \item Context output is used as extra input for arc inference.
\end{itemize}

\columnbreak

\textbf{Firmware deployment}
\begin{itemize}
  \item Exported \code{.h} headers are generated in \code{tinyml/model}.
  \item Firmware includes these generated headers during build.
  \item Rebuild and flash the firmware after updating active model headers.
  \item Keep the generation manifest and benchmark report with the exported model.
\end{itemize}

\textbf{Current active exported inputs}
\begin{itemize}
  \item Current RF header uses 13 inputs: 6 waveform/base features and 7 context inputs.
  \item Current context header uses 6 input features.
\end{itemize}
\end{multicols}

\begin{notebox}
The active exported RF header in \code{tinyml/model} uses 13 arc model inputs. An archived gen0 RF header used 12 inputs, but the current active header and manuscript workflow use the 13-input deployment structure.
\end{notebox}

\begin{layoutbox}
Use two side-by-side boxes: Arc Model and Context Model. Add arrows from Context Model to Arc Model.
\end{layoutbox}

\Spread{Spread 25}{Current Model Feature Reference}
\begin{tabularx}{\textwidth}{>{\bfseries}p{0.18\textwidth}X}
\toprule
Model & Current exported input order confirmed in active headers \\
\midrule
Arc RF base features & \code{max_low_current_run_ms}, \code{zero_dwell_ratio}, \code{low_current_ratio}, \code{abs_irms_zscore_vs_baseline}, \code{midband_residual_ratio}, \code{spectral_flux_midhf}. \\
Arc RF context inputs & \code{ctx_family_resistive_linear}, \code{ctx_family_inductive_motor}, \code{ctx_family_rectifier_smps}, \code{ctx_family_phase_angle_controlled}, \code{ctx_family_brush_universal_motor}, \code{ctx_family_other_mixed}, \code{context_family_confidence}. \\
Context model features & \code{delta_irms_abs}, \code{halfcycle_asymmetry}, \code{midband_residual_ratio}, \code{thd_i}, \code{spectral_flux_midhf}, \code{hf_energy_delta}. \\
Load family classes & \code{resistive_linear}, \code{inductive_motor}, \code{rectifier_smps}, \code{phase_angle_controlled}, \code{brush_universal_motor}, \code{other_mixed}. \\
\bottomrule
\end{tabularx}

\begin{notebox}
These lists are from the current exported headers in \code{tinyml/model}. If new headers are exported, update this page before printing the manual.
\end{notebox}

\begin{layoutbox}
Use this as a technical appendix-style spread. Keep it clean and table-based for developer reference.
\end{layoutbox}

\Spread{Spread 26}{Maintenance Guidelines}
\begin{multicols}{2}
\textbf{Before each use}
\begin{itemize}
  \item Check that the outlet is firm and undamaged.
  \item Check that the smart plug sits securely.
  \item Check the protected outlet and appliance plug.
  \item Keep the device dry and clear of clutter.
\end{itemize}

\textbf{After a trip}
\begin{itemize}
  \item Disconnect or turn off the appliance.
  \item Inspect for heat, smell, discoloration, looseness, or visible damage.
  \item Do not continue repeated tests without identifying the cause.
\end{itemize}

\columnbreak

\textbf{Storage}
\begin{itemize}
  \item Store in a dry indoor location.
  \item Do not place heavy items on the enclosure, plug, or button.
  \item Keep dust away from openings and connectors.
\end{itemize}

\textbf{Cleaning}
\begin{itemize}
  \item Turn the device OFF and unplug it.
  \item Use a dry cloth only.
  \item Do not spray liquid cleaner into the enclosure or socket.
\end{itemize}
\end{multicols}

\begin{warningbox}
Do not open the enclosure while the device is connected to mains power. Internal parts may carry dangerous voltage.
\end{warningbox}

\begin{layoutbox}
Use a four-card maintenance layout: Inspect, Clean, Store, After Fault.
\end{layoutbox}

\Spread{Spread 27}{Maintenance Schedule}
\begin{tabularx}{\textwidth}{>{\bfseries}p{0.18\textwidth}X X}
\toprule
When & Action & Purpose \\
\midrule
Before every use & Inspect outlet, plug fit, enclosure, protected socket, and appliance cord. & Prevent loose contact, overheating, and unsafe operation. \\
Before every use & Confirm the device is supported and not sagging. & Keep the plug tightly connected. \\
During use & Monitor OLED/PWA status and listen for chimes or alerts. & Catch warnings before continuing operation. \\
After any trip & Inspect load, socket, plug, wiring, and device temperature before clearing. & Avoid restarting into an unresolved fault. \\
Weekly during testing & Check for dust, discoloration, looseness, cracks, or button obstruction. & Maintain safe mechanical condition. \\
Before data collection & Confirm load label, phase, trial number, and logger settings. & Keep datasets organized and useful. \\
Before firmware update & Confirm firmware file, stable power, and correct device AP/browser workflow. & Reduce update failure risk. \\
\bottomrule
\end{tabularx}

\begin{layoutbox}
Use a clean schedule table. Color-code everyday, after-fault, and developer-only tasks.
\end{layoutbox}

\Spread{Spread 28}{Troubleshooting Guide}
\begin{tabularx}{\textwidth}{>{\bfseries}p{0.20\textwidth}X X}
\toprule
Problem & Possible cause & What to do \\
\midrule
Device does not power on & Main switch OFF, outlet has no power, loose plug, or damaged outlet. & Turn main switch ON. Try a known good outlet. Do not use loose or damaged outlets. \\
Appliance does not turn on & Relay is OFF by default, PWA controls are locked, fault is active, or load is not connected correctly. & Open PWA and turn relay ON if safe. Check fault status. Inspect before clearing any fault. \\
Cannot connect to Wi-Fi & Using 5 GHz network, wrong password, AP window closed, or weak signal. & Use 2.4 GHz Wi-Fi. Reset and try setup again. Confirm AP name and password. \\
Captive portal does not open & Browser or phone did not launch portal automatically. & Open browser and go to \code{192.168.4.1} while connected to the device AP. Turn off mobile data if needed. \\
PWA shows stale or offline & Device not connected to Wi-Fi/Firebase, browser cache, or network issue. & Refresh the PWA. Check Wi-Fi. Reset the device if needed. \\
Relay ON is blocked & Fault latch, web controls locked, startup, OTA, config portal, disconnected, or unsafe state. & Read the status message. Inspect the setup. Clear fault only after inspection. \\
Fault returns after clearing & The appliance, outlet, wiring, or load condition may still be unsafe. & Stop using the setup. Inspect and correct the cause before trying again. \\
Manual relay ON button does not work & No detectable load, not held long enough, fault condition, or hardware issue. & Connect the load first, hold about 5 seconds, wait for chime. Do not force or jam the button. \\
Logger stops but CSV still counts & Buffered samples are still uploading. & Wait until counting stops before viewing or downloading the CSV. \\
Firmware update page not found & Captive portal window, AP not open, mobile data active, or local update route unavailable. & Use a real browser at \code{192.168.4.1}, disable mobile data, and verify that the final firmware supports local Update Firmware. \\
\bottomrule
\end{tabularx}

\begin{layoutbox}
Use this as a full spread. Keep row height generous so Canva can place icons beside each problem.
\end{layoutbox}

\Spread{Spread 29}{Safety and Handling Checklist}
\begin{multicols}{2}
\textbf{Before powering ON}
\begin{itemize}
  \item Outlet is firm, dry, and undamaged.
  \item Device plug fits tightly.
  \item Smart plug is supported and not sagging.
  \item Appliance cord and plug are not damaged.
  \item Protected outlet is clear and dry.
  \item Manual relay ON button is not blocked or pressed.
\end{itemize}

\textbf{Before turning relay ON}
\begin{itemize}
  \item Correct load is connected.
  \item Load is within intended test conditions.
  \item No active fault is displayed.
  \item No unusual heat, smell, smoke, or sound is present.
\end{itemize}

\columnbreak

\textbf{After a warning or trip}
\begin{itemize}
  \item Turn off or disconnect the load.
  \item Inspect the socket, plug, wiring, and appliance.
  \item Let hot parts cool before touching nearby surfaces.
  \item Clear the fault only after inspection.
  \item Do not repeat operation if the same fault returns.
\end{itemize}

\textbf{For admin work}
\begin{itemize}
  \item Label the load, trial, and phase correctly.
  \item Keep raw and processed CSV files organized by generation.
  \item Do not update firmware on unstable power.
  \item Do not train or deploy models from poorly labeled data.
\end{itemize}
\end{multicols}

\begin{layoutbox}
Use checkboxes in Canva. This spread can also be printed as a lab checklist.
\end{layoutbox}

\Spread{Spread 30}{One-Page Quick Start Guide}
\begin{multicols}{3}
\textbf{1. Mount safely}
\begin{enumerate}
  \item Use a firm, undamaged outlet.
  \item Keep the device flat, flush, or supported.
  \item Do not let it hang from an extension cord or power strip.
\end{enumerate}

\textbf{2. Power the smart plug}
\begin{enumerate}
  \item Plug in the device.
  \item Turn the red main switch ON.
  \item Wait for startup.
\end{enumerate}

\textbf{3. Connect Wi-Fi}
\begin{enumerate}
  \item Connect to the \code{TinyML-Smart-Plug} AP.
  \item Open Configure Wi-Fi.
  \item Choose 2.4 GHz Wi-Fi.
  \item Enter credentials.
  \item Wait for chime and Wi-Fi connected status.
\end{enumerate}

\columnbreak

\textbf{4. Open dashboard}
\begin{enumerate}
  \item Open Chrome.
  \item Go to \code{tinyml-smart-plug.web.app}.
  \item Install the PWA if prompted.
\end{enumerate}

\textbf{5. Power the appliance}
\begin{enumerate}
  \item Plug the appliance into the protected outlet.
  \item Use the PWA relay rocker.
  \item Relay ON powers the appliance.
  \item Relay OFF turns the appliance off.
\end{enumerate}

\textbf{6. If a fault appears}
\begin{enumerate}
  \item Do not turn the relay back ON immediately.
  \item Inspect the load, socket, plug, and wiring.
  \item Click Clear Fault only after inspection.
\end{enumerate}

\columnbreak

\textbf{Offline relay ON}
\begin{enumerate}
  \item Connect the load first.
  \item Hold the red circular manual relay ON button for about 5 seconds.
  \item Release after the chime.
\end{enumerate}

\begin{warningbox}
The manual relay ON button cannot turn the relay OFF. It must never be stuck, blocked, or accidentally held.
\end{warningbox}

\begin{notebox}
Reset: red switch OFF, wait 10 seconds, red switch ON. Wi-Fi credentials remain saved.
\end{notebox}
\end{multicols}

\begin{layoutbox}
This is the printable quick-start insert. Use large numbered blocks and minimal text.
\end{layoutbox}

\Spread{Spread 31}{Screenshots and Photos to Capture}
\begin{tabularx}{\textwidth}{>{\bfseries}p{0.25\textwidth}X}
\toprule
Image & Use in booklet \\
\midrule
Front photo of the TinyML Smart Plug & Cover page and About page. \\
Device plugged into a stable wall outlet & Physical handling and mounting. \\
Incorrect sagging setup, staged safely with no power & Safety warning comparison. \\
Device on a flat power strip or extension setup & Alternative safe mounting. \\
Close-up of red main switch & Basic use and reset pages. \\
Close-up of manual relay ON button & Manual button instructions and warning. \\
OLED startup and Wi-Fi connected screen & Wi-Fi setup and offline protection pages. \\
Phone/laptop Wi-Fi AP list & Wi-Fi setup. \\
Captive portal Configure Wi-Fi screen & Wi-Fi setup and Change Wi-Fi. \\
PWA dashboard main page & Using the PWA and relay control. \\
Fault state and Clear Fault button & Fault behavior and clearing faults. \\
Top-right menu with User/Admin switch & Admin mode and Change Wi-Fi. \\
OTA admin section with green download arrow & Firmware update workflow. \\
Logger section & Data collection instructions. \\
CSV Plot Viewer with graph and series panel & CSV plotter/editor page. \\
Arc marker and segment editor controls & Arc labels and section labels page. \\
TinyML trainer folder or GUI screenshot & TinyML trainer workflow page. \\
\bottomrule
\end{tabularx}

\begin{layoutbox}
Use this as a production checklist while collecting Canva assets. Keep filenames consistent with the page numbers.
\end{layoutbox}

\Spread{Spread 32}{Source Confirmation Notes}
\begin{longtable}{>{\bfseries}p{0.30\textwidth}p{0.62\textwidth}}
\toprule
Item & Source-based decision used in this manual \\
\midrule
Wi-Fi AP name & Firmware constant uses \code{TinyML-Smart-Plug}; this manual uses that exact SSID. \\
AP timeout timing & Firmware uses about 15 seconds for boot AP wait without a client, about 45 seconds for normal portal setup after a client connects, about 45 seconds for manual AP wait, and about 4 minutes for manual portal setup. \\
OLED wording after Wi-Fi connect & Notification source draws \code{WIFI} and \code{CONNECTED}; this manual uses Wi-Fi connected status. \\
Manual OTA path & Main firmware supports HTTPS OTA through Firebase/PWA fields. Repository utility code includes a browser upload OTA route at \code{/update}; the manual local workflow is kept as developer-only. \\
Firmware release access & The manual states that firmware releases are developer-controlled, matching the user's operating instruction. \\
Relay hardware label & The manual uses the user-facing phrase relay output and protected socket. \\
Manual relay ON button behavior & Firmware confirms 5-second current-based physical relay ON confirmation; user hardware notes identify the red circular button. \\
Current active model input count & Active RF header uses 13 inputs. Archived gen0 RF header used 12, but the active deployment uses 13. \\
Firmware version & Source shows \code{v7.8.2-p-gen0-final}; this is printed as the firmware reference. \\
PWA menu wording & PWA source contains a top-right menu with User/Admin mode, Change WiFi, OTA, and notification controls. \\
Voltage protection wording & Firmware and manuscript include under-voltage and over-voltage states; this manual includes them as control-locking protection states. \\
\bottomrule
\end{longtable}

\begin{layoutbox}
Use this page as an appendix or internal source note. It explains why the manual uses the code-confirmed wording where earlier notes differed.
\end{layoutbox}

\Spread{Spread 33}{Contact and Support}
\begin{multicols}{2}
\textbf{For help}
\begin{itemize}
  \item Contact the thesis project developers or adviser.
  \item Provide the firmware version shown in the PWA.
  \item Describe the load used, outlet condition, and exact fault shown.
  \item Include screenshots of the PWA status and history entry if available.
\end{itemize}

\textbf{Before requesting support}
\begin{itemize}
  \item Stop using the device if there is smoke, smell, heat, sparking, or visible damage.
  \item Do not clear repeated faults without inspection.
  \item Do not open the enclosure while connected to mains power.
\end{itemize}

\columnbreak

\textbf{Support information to fill in}
\begin{tcolorbox}[colback=Panel,colframe=Line]
\textbf{Project team:} [FILL IN]\par
\textbf{Adviser:} [FILL IN]\par
\textbf{Email:} [FILL IN]\par
\textbf{Phone:} [FILL IN]\par
\textbf{Institution:} [FILL IN]\par
\textbf{Repo or release link:} [FILL IN]
\end{tcolorbox}
\end{multicols}

\begin{layoutbox}
Use a clean final page with contact fields, QR code to the PWA or repository if allowed, and a small product photo.
\end{layoutbox}

\Spread{Spread 34}{Glossary}
\begin{tabularx}{\textwidth}{>{\bfseries}p{0.22\textwidth}X}
\toprule
Term & Meaning \\
\midrule
Arc fault & An unsafe sparking or intermittent electrical condition that can occur through a damaged or loose connection. \\
CSV & Comma-separated values file used for logged sensor data and TinyML dataset preparation. \\
Estimated socket temperature & Firmware estimate of socket heating behavior based on sensor readings and model logic. \\
Firebase & Cloud service used by the project for live data, controls, event history, and logger uploads. \\
Firmware & Software running inside the ESP32-S3 smart plug controller. \\
Load & The appliance or test device plugged into the protected outlet. \\
OTA & Over-the-air firmware update process. \\
PWA & Progressive Web App. The dashboard can be opened in a browser and installed like an app. \\
Relay & Output switching path that energizes or disconnects the protected socket. \\
Sustained overload & Overload condition that remains long enough for the firmware to trip the relay. \\
TinyML & Machine learning model small enough to run locally on the microcontroller. \\
Trip & Protective action where the relay turns OFF because of a critical condition. \\
\bottomrule
\end{tabularx}

\begin{layoutbox}
Use as a final reference page. Add small icons only if space allows.
\end{layoutbox}

\end{document}
